遗憾的是, 边界项的估计相当细致, 因此将这个问题推迟到下一节. 先集中考虑高阶导数的内部估计. 例如, 设 $\Omega^{\prime}$ 为任意满足 $\Omega^{\prime}\subset\subset\Omega$ 的区域. (记号 $\Omega^{\prime}\subset\subset\Omega$ 表示 $\Omega^{\prime}$ 紧含于开集 $\Omega$; 即 $\overline{\Omega^{\prime}}$ 紧且 $\overline{\Omega^{\prime}}\subset\Omega$.) 现在选取截断函数 $\zeta\in\mathcal{D}(\Omega)$, 使 $0\leq\zeta\leq1$, 且在 $\Omega^{\prime}$ 上 $\zeta\equiv1$. 于是可以进行与前面十分相似的计算, 而没有边界项的干扰.

\[\begin{aligned}
\int_\Omega\zeta^2|\Delta u|^2\,\mathrm{d}\bm{x}
&=\sum_{i,j=1}^n\int_\Omega u_{x_ix_i}u_{x_jx_j}\zeta^2\,\mathrm{d}\bm{x}\\
&=-\sum_{i,j=1}^n\int_\Omega u_{x_i}u_{x_jx_ix_j}\zeta^2\,\mathrm{d}\bm{x}\\
&\quad-\sum_{i,j=1}^n\int_\Omega u_{x_i}u_{x_jx_j}2\zeta\zeta_{x_i}\,\mathrm{d}\bm{x}\\
&=\sum_{i,j=1}^n\int_\Omega u_{x_ix_j}u_{x_jx_i}\zeta^2\,\mathrm{d}\bm{x}\\
&\quad+\sum_{i,j=1}^n\int_\Omega u_{x_i}u_{x_jx_i}2\zeta\zeta_{x_j}-u_{x_i}u_{x_jx_j}2\zeta\zeta_{x_i}\,\mathrm{d}\bm{x}.
\end{aligned}\]

现在结合形如
\begin{equation}\label{eq:interior-cutoff-young-product-bound}
|u_{x_i}u_{x_ix_j}2\zeta\zeta_{x_j}|\leq\epsilon u_{x_ix_j}^2\zeta^2+\frac{1}{\epsilon}u_{x_i}^2\zeta_{x_j}^2
\end{equation}
的不等式和
\begin{equation}\label{eq:interior-cutoff-weighted-data-bound}
\int_\Omega|\Delta u|^2\zeta^2\,\mathrm{d}\bm{x}=\int_\Omega|f|^2\zeta^2\,\mathrm{d}\bm{x}\leq\int_\Omega|f|^2\,\mathrm{d}\bm{x}
\end{equation}
得到
\begin{equation}\label{eq:interior-cutoff-preabsorption-hessian-bound}
\int_\Omega|\nabla^2u|^2\zeta^2\,\mathrm{d}\bm{x}\leq\int_\Omega[|f|^2+\epsilon|\nabla^2u|^2\zeta^2+C(\epsilon)|\nabla u|^2|\nabla\zeta|^2]\,\mathrm{d}\bm{x}.
\end{equation}
再令 $\epsilon=1/2$, 并利用在 $\Omega^{\prime}$ 上 $\zeta\equiv1$, 得到
\begin{equation}\label{eq:interior-cutoff-local-hessian-apriori-estimate}
\int_{\Omega^{\prime}}|\nabla^2u|^2\,\mathrm{d}\bm{x}\leq\int_\Omega\zeta^2|\nabla^2u|^2\,\mathrm{d}\bm{x}\leq C\left(\|f\|_2^2+\|\nabla u\|_2^2\right).
\end{equation}


因此, 对任意 $\Omega^{\prime}\subset\subset\Omega$, 都得到了用数据 $f$ 的 $L^2(\Omega)$ 和 $u$ 的 $H^1(\Omega)$ 范数控制解 $u$ 的 $H^2(\Omega^{\prime})$ 范数的估计. 当然, 对上述计算的一个主要异议是, 为进行其中的分部积分, 必须对解 $u$ 的光滑性作未经论证的假设. 在下面这些计算的严格版本中, 将用涉及差商的类似技巧替代这些操作. 由于差商技巧在本节十分重要, 下面对这一主题作简要讨论.
